% =============================================================================
% Chapter 17: Data Streams - Λυμένες Ασκήσεις
% Data Mining Study Guide - NTUA
% =============================================================================

\chapter{Data Streams - Λυμένες Ασκήσεις}
\label{ch:data-streams-exercises}

Αυτό το κεφάλαιο περιέχει λυμένες ασκήσεις για τους κύριους αλγορίθμους data streams:
\begin{itemize}
    \item \textbf{Bloom Filters}: Membership queries και optimal sizing (Ασκήσεις 6.1, 6.2, 6.11)
    \item \textbf{Count-Min Sketch}: Frequency estimation και heavy hitters (Ασκήσεις 6.3--6.5)
    \item \textbf{Flajolet-Martin}: Distinct element counting (Ασκήσεις 6.7, 6.8)
    \item \textbf{Reservoir Sampling}: Uniform και biased sampling (Ασκήσεις 6.9, 6.10)
    \item \textbf{Algorithm Selection}: Επιλογή κατάλληλης δομής (Ασκήσεις 6.6, 6.12)
    \item \textbf{STREAM Algorithm}: Clustering με Cluster Features (Άσκηση 6.13)
    \item \textbf{Lossy Counting}: Frequent items detection (Άσκηση 6.14)
\end{itemize}

% -----------------------------------------------------------------------------
% Exercise 6.1
% -----------------------------------------------------------------------------
\section{Άσκηση 6.1: Bloom Filter Insert \& Lookup}

\begin{formulabox}[Εκφώνηση]
Ένα Bloom Filter έχει μέγεθος $m = 10$ bits και χρησιμοποιεί 2 hash functions:
\begin{itemize}
    \item $h_1(x) = x \mod 10$
    \item $h_2(x) = (2x + 3) \mod 10$
\end{itemize}

\begin{enumerate}
    \item Εισάγετε τα στοιχεία: 15, 23, 36
    \item Ελέγξτε αν υπάρχουν: 23, 42
\end{enumerate}
\end{formulabox}

\subsection*{Λύση}

\textbf{Βήμα 1: Αρχικοποίηση}

Bit array: $[0, 0, 0, 0, 0, 0, 0, 0, 0, 0]$

\textbf{Βήμα 2: Insert 15}

\begin{align*}
h_1(15) &= 15 \mod 10 = 5 \\
h_2(15) &= (30 + 3) \mod 10 = 3
\end{align*}

Set bits 3, 5: $[0, 0, 0, \mathbf{1}, 0, \mathbf{1}, 0, 0, 0, 0]$

\textbf{Βήμα 3: Insert 23}

\begin{align*}
h_1(23) &= 23 \mod 10 = 3 \\
h_2(23) &= (46 + 3) \mod 10 = 9
\end{align*}

Set bits 3, 9: $[0, 0, 0, \mathbf{1}, 0, 1, 0, 0, 0, \mathbf{1}]$

\textbf{Βήμα 4: Insert 36}

\begin{align*}
h_1(36) &= 36 \mod 10 = 6 \\
h_2(36) &= (72 + 3) \mod 10 = 5
\end{align*}

Set bits 5, 6: $[0, 0, 0, 1, 0, \mathbf{1}, \mathbf{1}, 0, 0, 1]$

\textbf{Final bit array:} $[0, 0, 0, 1, 0, 1, 1, 0, 0, 1]$

\textbf{Lookup 23:}
\begin{itemize}
    \item $h_1(23) = 3$ $\to$ bit[3] = 1 $\checkmark$
    \item $h_2(23) = 9$ $\to$ bit[9] = 1 $\checkmark$
\end{itemize}
Αποτέλεσμα: \textbf{Probably in set}

\textbf{Lookup 42:}
\begin{align*}
h_1(42) &= 42 \mod 10 = 2 \\
h_2(42) &= (84 + 3) \mod 10 = 7
\end{align*}
\begin{itemize}
    \item bit[2] = 0 $\times$
\end{itemize}
Αποτέλεσμα: \textbf{Definitely not in set}

% -----------------------------------------------------------------------------
% Exercise 6.2
% -----------------------------------------------------------------------------
\section{Άσκηση 6.2: Bloom Filter False Positive}

\begin{formulabox}[Εκφώνηση]
Με το ίδιο Bloom Filter από την προηγούμενη άσκηση, ελέγξτε αν το 13 είναι
false positive.
\end{formulabox}

\subsection*{Λύση}

Bit array: $[0, 0, 0, 1, 0, 1, 1, 0, 0, 1]$

\textbf{Lookup 13:}
\begin{align*}
h_1(13) &= 13 \mod 10 = 3 \\
h_2(13) &= (26 + 3) \mod 10 = 9
\end{align*}

\begin{itemize}
    \item bit[3] = 1 $\checkmark$
    \item bit[9] = 1 $\checkmark$
\end{itemize}

Αποτέλεσμα: «Probably in set»

\textbf{Αλλά:} Το 13 \textbf{δεν} εισήχθη ποτέ! Αυτό είναι \textbf{FALSE POSITIVE}.

\begin{warningbox}[False Positive Rate]
Η πιθανότητα false positive είναι περίπου:
$$P_{fp} \approx \left(1 - e^{-kn/m}\right)^k$$
όπου $k$ = hash functions, $n$ = inserted elements, $m$ = bits.

Για $k=2$, $n=3$, $m=10$:
$$P_{fp} \approx (1 - e^{-0.6})^2 \approx (0.45)^2 \approx 0.20 = 20\%$$
\end{warningbox}

% -----------------------------------------------------------------------------
% Exercise 6.3
% -----------------------------------------------------------------------------
\section{Άσκηση 6.3: Count-Min Sketch Query}

\begin{formulabox}[Εκφώνηση]
Ένα Count-Min Sketch έχει $d = 2$ hash functions και $w = 5$ counters ανά row.
Μετά από updates, ο πίνακας είναι:

\begin{center}
\begin{tabular}{c|ccccc}
 & 0 & 1 & 2 & 3 & 4 \\
\hline
$h_1$ & 3 & 7 & 2 & 5 & 1 \\
$h_2$ & 4 & 2 & 8 & 3 & 6 \\
\end{tabular}
\end{center}

Αν $h_1(\text{``apple''}) = 1$ και $h_2(\text{``apple''}) = 3$, ποια είναι η εκτίμηση
για το count του ``apple''?
\end{formulabox}

\subsection*{Λύση}

\textbf{Count-Min Sketch Query:}

$$\hat{f}(x) = \min_{i=1}^{d} \text{CMS}[i][h_i(x)]$$

\textbf{Υπολογισμός:}
\begin{align*}
\text{CMS}[1][h_1(\text{``apple''})] &= \text{CMS}[1][1] = 7 \\
\text{CMS}[2][h_2(\text{``apple''})] &= \text{CMS}[2][3] = 3
\end{align*}

$$\hat{f}(\text{``apple''}) = \min(7, 3) = 3$$

\begin{infobox}[Γιατί Minimum?]
Τα counts μπορούν να αυξηθούν λόγω collisions με άλλα στοιχεία.
Το minimum δίνει την καλύτερη (μικρότερη) εκτίμηση, που είναι εγγυημένα $\geq$ actual count.
\end{infobox}

% -----------------------------------------------------------------------------
% Exercise 6.4
% -----------------------------------------------------------------------------
\section{Άσκηση 6.4: Count-Min Sketch Update}

\begin{formulabox}[Εκφώνηση]
Χρησιμοποιώντας τον ίδιο CMS, εκτελέστε: Update(``banana'', +2)

Δίνεται: $h_1(\text{``banana''}) = 4$, $h_2(\text{``banana''}) = 2$
\end{formulabox}

\subsection*{Λύση}

\textbf{Update operation:}

Για κάθε row $i$: $\text{CMS}[i][h_i(x)] \mathrel{+}= c$

\textbf{Πριν:}
\begin{center}
\begin{tabular}{c|ccccc}
 & 0 & 1 & 2 & 3 & 4 \\
\hline
$h_1$ & 3 & 7 & 2 & 5 & 1 \\
$h_2$ & 4 & 2 & 8 & 3 & 6 \\
\end{tabular}
\end{center}

\textbf{Update:}
\begin{align*}
\text{CMS}[1][4] &= 1 + 2 = 3 \\
\text{CMS}[2][2] &= 8 + 2 = 10
\end{align*}

\textbf{Μετά:}
\begin{center}
\begin{tabular}{c|ccccc}
 & 0 & 1 & 2 & 3 & 4 \\
\hline
$h_1$ & 3 & 7 & 2 & 5 & \textbf{3} \\
$h_2$ & 4 & 2 & \textbf{10} & 3 & 6 \\
\end{tabular}
\end{center}

% -----------------------------------------------------------------------------
% Exercise 6.5
% -----------------------------------------------------------------------------
\section{Άσκηση 6.5: Error Bounds}

\begin{formulabox}[Εκφώνηση]
Σχεδιάστε Count-Min Sketch με error $\varepsilon = 0.01$ και failure probability
$\delta = 0.05$. Υπολογίστε τις διαστάσεις.
\end{formulabox}

\subsection*{Λύση}

\textbf{Τύποι:}

$$w = \lceil e / \varepsilon \rceil \approx \lceil 2.718 / 0.01 \rceil = 272$$

$$d = \lceil \ln(1/\delta) \rceil = \lceil \ln(20) \rceil = \lceil 3.0 \rceil = 3$$

\textbf{Εγγύηση:}

Με πιθανότητα τουλάχιστον $1 - \delta = 95\%$:
$$\hat{f}(x) \leq f(x) + \varepsilon \cdot N$$

όπου $N$ = συνολικό count όλων των updates.

\textbf{Απαιτούμενη μνήμη:}
$$\text{Memory} = w \times d \times \text{counter\_size} = 272 \times 3 \times 4 = 3264 \text{ bytes}$$

% -----------------------------------------------------------------------------
% Exercise 6.6
% -----------------------------------------------------------------------------
\section{Άσκηση 6.6: Bloom Filter vs Count-Min Sketch}

\begin{formulabox}[Εκφώνηση]
Για κάθε σενάριο, επιλέξτε την κατάλληλη δομή (Bloom Filter ή Count-Min Sketch):

\begin{enumerate}
    \item Έλεγχος αν ένα URL έχει επισκεφθεί πριν
    \item Καταμέτρηση επισκέψεων ανά URL
    \item Spam filter: έλεγχος αν email είναι σε blacklist
    \item Εύρεση heavy hitters σε network traffic
\end{enumerate}
\end{formulabox}

\subsection*{Λύση}

\begin{center}
\begin{tabular}{clp{6cm}}
\toprule
\textbf{\#} & \textbf{Δομή} & \textbf{Αιτιολόγηση} \\
\midrule
1 & \textbf{Bloom Filter} & Membership test μόνο (yes/no) \\
2 & \textbf{Count-Min Sketch} & Χρειάζεται καταμέτρηση, όχι μόνο membership \\
3 & \textbf{Bloom Filter} & Membership test σε blacklist \\
4 & \textbf{Count-Min Sketch} & Heavy hitters = υψηλά counts \\
\bottomrule
\end{tabular}
\end{center}

\begin{infobox}[Σύγκριση]
\begin{itemize}
    \item \textbf{Bloom Filter}: Membership queries ($\in$?), no deletions, no counts
    \item \textbf{Count-Min Sketch}: Frequency queries (count?), supports negative updates
\end{itemize}
\end{infobox}

% -----------------------------------------------------------------------------
% Exercise 6.7
% -----------------------------------------------------------------------------
\section{Άσκηση 6.7: Flajolet-Martin Basic Calculation}

\begin{formulabox}[Εκφώνηση]
Θέλουμε να εκτιμήσουμε τον αριθμό των διακριτών στοιχείων σε ένα data stream
χρησιμοποιώντας τον αλγόριθμο Flajolet-Martin.

Δίνεται η hash function $h(x)$ που επιστρέφει 8-bit τιμές. Τα στοιχεία του stream
και οι hash τιμές τους (σε binary) είναι:

\begin{center}
\begin{tabular}{ccc}
\toprule
\textbf{Element} & \textbf{h(x)} & \textbf{Binary} \\
\midrule
A & 44 & 00101100 \\
B & 56 & 00111000 \\
C & 16 & 00010000 \\
A & 44 & 00101100 \\
D & 72 & 01001000 \\
B & 56 & 00111000 \\
E & 32 & 00100000 \\
\bottomrule
\end{tabular}
\end{center}

\begin{enumerate}
    \item Υπολογίστε το $R$ (trailing zeros) για κάθε στοιχείο
    \item Βρείτε το $R_{max}$
    \item Εκτιμήστε τον αριθμό διακριτών στοιχείων
    \item Ποιος είναι ο πραγματικός αριθμός διακριτών στοιχείων;
\end{enumerate}
\end{formulabox}

\subsection*{Λύση}

\textbf{Βήμα 1: Υπολογισμός trailing zeros ($R$)}

Το $R$ είναι ο αριθμός των trailing zeros στη binary hash τιμή:

\begin{center}
\begin{tabular}{cccc}
\toprule
\textbf{Element} & \textbf{Binary} & \textbf{Trailing Zeros} & \textbf{R} \\
\midrule
A & 0010110\underline{0} & 0 & 2 \\
B & 0011100\underline{0} & 000 & 3 \\
C & 0001000\underline{0} & 0000 & 4 \\
D & 0100100\underline{0} & 000 & 3 \\
E & 0010000\underline{0} & 00000 & 5 \\
\bottomrule
\end{tabular}
\end{center}

\textbf{Βήμα 2: $R_{max}$}

$$R_{max} = \max(2, 3, 4, 3, 5) = 5$$

\textbf{Βήμα 3: Εκτίμηση διακριτών στοιχείων}

$$\text{Estimate} = 2^{R_{max}} = 2^5 = 32$$

\textbf{Βήμα 4: Πραγματικός αριθμός}

Τα διακριτά στοιχεία είναι: $\{A, B, C, D, E\}$

$$\text{Actual distinct} = 5$$

\begin{warningbox}[FM Estimation Error]
Η εκτίμηση (32) είναι πολύ μεγαλύτερη από την πραγματική τιμή (5).
Αυτό είναι χαρακτηριστικό του FM algorithm με μία hash function.
Η ακρίβεια βελτιώνεται σημαντικά με πολλαπλές hash functions (βλ. Άσκηση 6.8).
\end{warningbox}

% -----------------------------------------------------------------------------
% Exercise 6.8
% -----------------------------------------------------------------------------
\section{Άσκηση 6.8: Flajolet-Martin με Πολλαπλές Hash Functions}

\begin{formulabox}[Εκφώνηση]
Για να βελτιώσουμε την ακρίβεια του FM algorithm, χρησιμοποιούμε 6 hash functions
χωρισμένες σε 2 groups των 3. Οι $R_{max}$ τιμές για κάθε hash function είναι:

\begin{center}
\begin{tabular}{cc}
\toprule
\textbf{Group 1} & \textbf{Group 2} \\
\midrule
$R^{(1)}_{max} = 3$ & $R^{(4)}_{max} = 4$ \\
$R^{(2)}_{max} = 5$ & $R^{(5)}_{max} = 2$ \\
$R^{(3)}_{max} = 4$ & $R^{(6)}_{max} = 3$ \\
\bottomrule
\end{tabular}
\end{center}

Υπολογίστε την τελική εκτίμηση χρησιμοποιώντας τη μέθοδο median-of-averages.
\end{formulabox}

\subsection*{Λύση}

\textbf{Βήμα 1: Υπολογισμός εκτιμήσεων ανά hash}

\begin{align*}
\text{Est}_1 &= 2^3 = 8 & \text{Est}_4 &= 2^4 = 16 \\
\text{Est}_2 &= 2^5 = 32 & \text{Est}_5 &= 2^2 = 4 \\
\text{Est}_3 &= 2^4 = 16 & \text{Est}_6 &= 2^3 = 8
\end{align*}

\textbf{Βήμα 2: Μέσος όρος ανά group}

\begin{align*}
\text{Avg}_{\text{Group 1}} &= \frac{8 + 32 + 16}{3} = \frac{56}{3} \approx 18.67 \\
\text{Avg}_{\text{Group 2}} &= \frac{16 + 4 + 8}{3} = \frac{28}{3} \approx 9.33
\end{align*}

\textbf{Βήμα 3: Median των averages}

$$\text{Estimate} = \text{median}(18.67, 9.33) = \frac{18.67 + 9.33}{2} = 14$$

\begin{infobox}[Median-of-Averages Method]
\begin{itemize}
    \item \textbf{Average}: Μειώνει τη διακύμανση (variance)
    \item \textbf{Median}: Προστατεύει από extreme τιμές (outliers)
    \item Συνήθως: $m$ groups × $k$ hash functions = $m \cdot k$ συνολικά
    \item Accuracy βελτιώνεται ως $O(1/\sqrt{mk})$
\end{itemize}
\end{infobox}

% -----------------------------------------------------------------------------
% Exercise 6.9
% -----------------------------------------------------------------------------
\section{Άσκηση 6.9: Reservoir Sampling Basic}

\begin{formulabox}[Εκφώνηση]
Έχουμε ένα data stream με στοιχεία που έρχονται ένα-ένα:
$$\text{Stream}: A, B, C, D, E, F, G, H$$

Θέλουμε να διατηρήσουμε uniform random sample μεγέθους $k = 3$ χρησιμοποιώντας
Reservoir Sampling. Δίνονται οι τυχαίοι αριθμοί $r \in [0,1)$ για κάθε βήμα:

\begin{center}
\begin{tabular}{cc}
\toprule
\textbf{Step (n)} & \textbf{Random r} \\
\midrule
4 (D) & 0.65 \\
5 (E) & 0.45 \\
6 (F) & 0.33 \\
7 (G) & 0.12 \\
8 (H) & 0.80 \\
\bottomrule
\end{tabular}
\end{center}

Εκτελέστε τον αλγόριθμο και δείξτε την κατάσταση του reservoir σε κάθε βήμα.
\end{formulabox}

\subsection*{Λύση}

\textbf{Reservoir Sampling Algorithm:}
\begin{enumerate}
    \item Για $n \leq k$: Βάλε το στοιχείο στο reservoir
    \item Για $n > k$: Με πιθανότητα $k/n$, αντικατέστησε τυχαίο στοιχείο
\end{enumerate}

Replacement γίνεται αν $r < k/n$. Αν ναι, αντικαθιστούμε position $\lfloor r \cdot n \rfloor \mod k$.

\textbf{Εκτέλεση:}

\begin{center}
\begin{tabular}{cccccl}
\toprule
\textbf{n} & \textbf{Element} & \textbf{k/n} & \textbf{r} & \textbf{r < k/n?} & \textbf{Reservoir} \\
\midrule
1 & A & -- & -- & -- & [A, -, -] \\
2 & B & -- & -- & -- & [A, B, -] \\
3 & C & -- & -- & -- & [A, B, C] \\
4 & D & 3/4=0.75 & 0.65 & Yes & [A, B, \textbf{D}] \\
5 & E & 3/5=0.60 & 0.45 & Yes & [A, \textbf{E}, D] \\
6 & F & 3/6=0.50 & 0.33 & Yes & [\textbf{F}, E, D] \\
7 & G & 3/7=0.43 & 0.12 & Yes & [F, E, \textbf{G}] \\
8 & H & 3/8=0.375 & 0.80 & No & [F, E, G] \\
\bottomrule
\end{tabular}
\end{center}

\textbf{Position υπολογισμός} (όταν $r < k/n$):
\begin{itemize}
    \item $n=4$: $\lfloor 0.65 \cdot 4 \rfloor \mod 3 = 2 \mod 3 = 2$ $\to$ αντικατάσταση pos 2
    \item $n=5$: $\lfloor 0.45 \cdot 5 \rfloor \mod 3 = 2 \mod 3 = 2$ $\to$ αντικατάσταση pos 1
    \item $n=6$: $\lfloor 0.33 \cdot 6 \rfloor \mod 3 = 1 \mod 3 = 1$ $\to$ αντικατάσταση pos 0
    \item $n=7$: $\lfloor 0.12 \cdot 7 \rfloor \mod 3 = 0 \mod 3 = 0$ $\to$ αντικατάσταση pos 2
\end{itemize}

\textbf{Final Reservoir:} $\{F, E, G\}$

\begin{infobox}[Uniform Sampling Property]
Κάθε στοιχείο του stream έχει ίση πιθανότητα $k/n$ να βρίσκεται στο τελικό sample.
Η απόδειξη γίνεται με induction: αν μετά από $n-1$ στοιχεία κάθε ένα έχει
πιθανότητα $k/(n-1)$, τότε μετά το $n$-οστό:
$$P(\text{element in sample}) = \frac{k}{n-1} \cdot \left(1 - \frac{k}{n} \cdot \frac{1}{k}\right) = \frac{k}{n-1} \cdot \frac{n-1}{n} = \frac{k}{n}$$
\end{infobox}

% -----------------------------------------------------------------------------
% Exercise 6.10
% -----------------------------------------------------------------------------
\section{Άσκηση 6.10: Reservoir Sampling με Exponential Bias}

\begin{formulabox}[Εκφώνηση]
Σε περιβάλλοντα με concept drift, θέλουμε τα πιο πρόσφατα στοιχεία να έχουν
μεγαλύτερη πιθανότητα να επιλεγούν. Χρησιμοποιούμε exponential bias function:

$$f(r, n) = e^{-\lambda(n - r)}$$

όπου $r$ = position του στοιχείου στο stream, $n$ = τρέχον μήκος stream,
$\lambda$ = decay rate.

Δίνεται: $k = 2$ (reservoir size), $\lambda = 0.5$

Stream: $X_1, X_2, X_3, X_4, X_5$

\begin{enumerate}
    \item Υπολογίστε το bias weight $f(r, n)$ για κάθε στοιχείο όταν $n = 5$
    \item Υπολογίστε την πιθανότητα insertion για το νέο στοιχείο $X_5$
    \item Υπολογίστε το $F(n) = \sum_{r=1}^{n} f(r, n)$ και την πιθανότητα replacement
\end{enumerate}
\end{formulabox}

\subsection*{Λύση}

\textbf{Βήμα 1: Bias weights για $n = 5$}

\begin{center}
\begin{tabular}{cccc}
\toprule
\textbf{Element} & \textbf{r} & \textbf{n - r} & \textbf{$f(r,5) = e^{-0.5(5-r)}$} \\
\midrule
$X_1$ & 1 & 4 & $e^{-2.0} = 0.135$ \\
$X_2$ & 2 & 3 & $e^{-1.5} = 0.223$ \\
$X_3$ & 3 & 2 & $e^{-1.0} = 0.368$ \\
$X_4$ & 4 & 1 & $e^{-0.5} = 0.607$ \\
$X_5$ & 5 & 0 & $e^{0} = 1.000$ \\
\bottomrule
\end{tabular}
\end{center}

\textbf{Βήμα 2: Insertion probability}

Το νέο στοιχείο $X_n$ (με $f(n,n) = 1$) εισάγεται με πιθανότητα:
$$P(\text{insert } X_n) = \lambda k = 0.5 \times 2 = 1.0$$

Αυτό σημαίνει ότι κάθε νέο στοιχείο εισάγεται πάντα (για αυτό το $\lambda k$).

\textbf{Βήμα 3: Total weight και replacement probability}

$$F(5) = \sum_{r=1}^{5} e^{-0.5(5-r)} = 0.135 + 0.223 + 0.368 + 0.607 + 1.0 = 2.333$$

Η πιθανότητα να αντικατασταθεί συγκεκριμένο παλιό στοιχείο $X_r$:
$$P(\text{replace } X_r) = \frac{f(r, n)}{F(n)} = \frac{e^{-\lambda(n-r)}}{\sum_{j=1}^{n} e^{-\lambda(n-j)}}$$

\begin{warningbox}[Exponential Decay Effect]
\begin{itemize}
    \item $X_1$ (oldest): weight = 0.135, relative prob = $0.135/2.333 = 5.8\%$
    \item $X_5$ (newest): weight = 1.000, relative prob = $1.0/2.333 = 42.9\%$
\end{itemize}
Τα πιο πρόσφατα στοιχεία έχουν $\sim$7x μεγαλύτερη πιθανότητα να βρίσκονται στο sample!
\end{warningbox}

% -----------------------------------------------------------------------------
% Exercise 6.11
% -----------------------------------------------------------------------------
\section{Άσκηση 6.11: Bloom Filter Optimal Sizing}

\begin{formulabox}[Εκφώνηση]
Θέλουμε να σχεδιάσουμε ένα Bloom Filter για να αποθηκεύσουμε $n = 10{,}000$
email addresses (για spam detection) με target false positive rate $F = 1\%$.

\begin{enumerate}
    \item Υπολογίστε το βέλτιστο μέγεθος $m$ (bits)
    \item Υπολογίστε τον βέλτιστο αριθμό hash functions $k$
    \item Επαληθεύστε ότι επιτυγχάνεται το target false positive rate
    \item Πόση μνήμη απαιτείται;
\end{enumerate}
\end{formulabox}

\subsection*{Λύση}

\textbf{Βήμα 1: Optimal filter size $m$}

Ο τύπος για optimal $m$ δίνεται από:
$$m = -\frac{n \ln F}{(\ln 2)^2}$$

\begin{align*}
m &= -\frac{10{,}000 \times \ln(0.01)}{(\ln 2)^2} \\
&= -\frac{10{,}000 \times (-4.605)}{0.693^2} \\
&= \frac{46{,}050}{0.480} \\
&\approx 95{,}937 \text{ bits}
\end{align*}

Στρογγυλοποιούμε: $m = 96{,}000$ bits

\textbf{Βήμα 2: Optimal number of hash functions $k$}

$$k = \frac{m}{n} \ln 2 = \frac{96{,}000}{10{,}000} \times 0.693 \approx 6.65$$

Στρογγυλοποιούμε: $k = 7$ hash functions

\textbf{Βήμα 3: Επαλήθευση false positive rate}

$$P_{fp} = \left(1 - e^{-kn/m}\right)^k = \left(1 - e^{-7 \times 10{,}000 / 96{,}000}\right)^7$$

\begin{align*}
P_{fp} &= \left(1 - e^{-0.729}\right)^7 \\
&= (1 - 0.482)^7 \\
&= (0.518)^7 \\
&\approx 0.0097 \approx 0.97\%
\end{align*}

\textbf{Βήμα 4: Memory requirements}

$$\text{Memory} = \frac{m}{8} = \frac{96{,}000}{8} = 12{,}000 \text{ bytes} = 11.7 \text{ KB}$$

\begin{infobox}[Bits per Element]
Ο λόγος $m/n$ (bits per element) είναι σταθερός για δεδομένο $F$:
$$\frac{m}{n} = -\frac{\ln F}{(\ln 2)^2} \approx 1.44 \log_2(1/F)$$

Για $F = 1\%$: $m/n \approx 9.6$ bits/element
\end{infobox}

% -----------------------------------------------------------------------------
% Exercise 6.12
% -----------------------------------------------------------------------------
\section{Άσκηση 6.12: Επιλογή Αλγορίθμου για Data Streams}

\begin{formulabox}[Εκφώνηση]
Για κάθε σενάριο, επιλέξτε τον καταλληλότερο αλγόριθμο/δομή δεδομένων και
αιτιολογήστε την επιλογή σας.

\textbf{Διαθέσιμες επιλογές:}
\begin{itemize}
    \item Bloom Filter
    \item Count-Min Sketch
    \item Flajolet-Martin (HyperLogLog)
    \item Reservoir Sampling
\end{itemize}

\textbf{Σενάρια:}
\begin{enumerate}
    \item Υπολογισμός του αριθμού unique visitors σε website (εκατομμύρια επισκέψεις)
    \item Διατήρηση representative sample από sensor data για statistical analysis
    \item Έλεγχος αν ένα password έχει διαρρεύσει (leaked passwords database)
    \item Εύρεση των top-10 trending hashtags σε social media stream
    \item Ανίχνευση DDoS attack: counting requests per IP address
    \item Random sampling από log files για debugging (χωρίς να φορτώσουμε όλο το αρχείο)
\end{enumerate}
\end{formulabox}

\subsection*{Λύση}

\begin{center}
\renewcommand{\arraystretch}{1.4}
\begin{tabular}{clp{7cm}}
\toprule
\textbf{\#} & \textbf{Αλγόριθμος} & \textbf{Αιτιολόγηση} \\
\midrule
1 & \textbf{Flajolet-Martin} & Distinct count estimation με $O(\log \log n)$ space. Ideal για cardinality estimation σε μεγάλα datasets. \\
\midrule
2 & \textbf{Reservoir Sampling} & Διατηρεί uniform random sample σταθερού μεγέθους, κατάλληλο για statistical analysis. \\
\midrule
3 & \textbf{Bloom Filter} & Membership test (``έχει διαρρεύσει;'') με χαμηλό false positive rate. No false negatives = ασφαλές για passwords. \\
\midrule
4 & \textbf{Count-Min Sketch} & Heavy hitters detection: track frequency counts και βρες τα στοιχεία με highest counts. \\
\midrule
5 & \textbf{Count-Min Sketch} & Frequency estimation per IP. Threshold στα counts για anomaly detection. \\
\midrule
6 & \textbf{Reservoir Sampling} & Single-pass sampling χωρίς να ξέρουμε το μέγεθος εκ των προτέρων. Perfect για log files. \\
\bottomrule
\end{tabular}
\end{center}

\begin{infobox}[Decision Tree για Algorithm Selection]
\begin{enumerate}
    \item \textbf{Τι query θέλω;}
    \begin{itemize}
        \item Membership ($x \in S$?) $\to$ Bloom Filter
        \item Frequency (πόσες φορές $x$?) $\to$ Count-Min Sketch
        \item Cardinality ($|S|$?) $\to$ Flajolet-Martin / HyperLogLog
        \item Sample $\to$ Reservoir Sampling
    \end{itemize}
    \item \textbf{Space constraint?}
    \begin{itemize}
        \item Bloom: $O(n)$ bits
        \item CMS: $O(1/\varepsilon \cdot \log(1/\delta))$
        \item FM: $O(\log \log n)$
        \item Reservoir: $O(k)$
    \end{itemize}
    \item \textbf{Error tolerance?}
    \begin{itemize}
        \item Bloom: False positives OK, no false negatives
        \item CMS: Overestimates OK, never underestimates
        \item FM: Approximate count OK
    \end{itemize}
\end{enumerate}
\end{infobox}

% -----------------------------------------------------------------------------
% Exercise 6.13
% -----------------------------------------------------------------------------
\section{Άσκηση 6.13: STREAM Algorithm (Clustering)}

\begin{formulabox}[Εκφώνηση]
Ο αλγόριθμος STREAM είναι μια προσαρμογή του k-means για data streams.
Χρησιμοποιεί τη δομή BIRCH CF-tree με Cluster Features (CF).

Δίνονται 3 micro-clusters με Cluster Features:
\begin{center}
\begin{tabular}{cccc}
\toprule
\textbf{CF} & \textbf{N} (points) & \textbf{LS} (linear sum) & \textbf{SS} (squared sum) \\
\midrule
$CF_1$ & 100 & $(200, 300)$ & $(500, 1000)$ \\
$CF_2$ & 50 & $(150, 100)$ & $(550, 250)$ \\
$CF_3$ & 75 & $(300, 225)$ & $(1350, 750)$ \\
\bottomrule
\end{tabular}
\end{center}

\begin{enumerate}
    \item Υπολογίστε το centroid κάθε micro-cluster
    \item Υπολογίστε το radius (RMSD) του $CF_1$
    \item Αν συγχωνεύσουμε $CF_1$ και $CF_2$, ποιο είναι το νέο CF;
\end{enumerate}
\end{formulabox}

\subsection*{Λύση}

\textbf{Cluster Feature (CF) ορισμός:}

Για cluster $C$ με $N$ σημεία $\{x_1, ..., x_N\}$:
$$CF = (N, \vec{LS}, \vec{SS})$$
όπου:
\begin{itemize}
    \item $\vec{LS} = \sum_{i=1}^{N} x_i$ (linear sum)
    \item $\vec{SS} = \sum_{i=1}^{N} x_i^2$ (squared sum, element-wise)
\end{itemize}

\textbf{(1) Υπολογισμός Centroids:}

$$\text{Centroid} = \frac{\vec{LS}}{N}$$

\begin{align*}
C_1 &= \frac{(200, 300)}{100} = (2.0, 3.0) \\
C_2 &= \frac{(150, 100)}{50} = (3.0, 2.0) \\
C_3 &= \frac{(300, 225)}{75} = (4.0, 3.0)
\end{align*}

\textbf{(2) Υπολογισμός Radius (RMSD) για $CF_1$:}

$$R = \sqrt{\frac{\vec{SS}}{N} - \left(\frac{\vec{LS}}{N}\right)^2}$$

Για κάθε διάσταση:
\begin{align*}
R_x &= \sqrt{\frac{500}{100} - \left(\frac{200}{100}\right)^2} = \sqrt{5 - 4} = 1.0 \\
R_y &= \sqrt{\frac{1000}{100} - \left(\frac{300}{100}\right)^2} = \sqrt{10 - 9} = 1.0
\end{align*}

$$R_1 = \sqrt{R_x^2 + R_y^2} = \sqrt{1 + 1} = \sqrt{2} \approx 1.41$$

\textbf{(3) Συγχώνευση $CF_1$ και $CF_2$:}

Η additive property των CFs επιτρέπει απλή πρόσθεση:
$$CF_{merged} = CF_1 + CF_2 = (N_1 + N_2, \vec{LS}_1 + \vec{LS}_2, \vec{SS}_1 + \vec{SS}_2)$$

\begin{align*}
N_{merged} &= 100 + 50 = 150 \\
\vec{LS}_{merged} &= (200 + 150, 300 + 100) = (350, 400) \\
\vec{SS}_{merged} &= (500 + 550, 1000 + 250) = (1050, 1250)
\end{align*}

$$CF_{merged} = (150, (350, 400), (1050, 1250))$$

\begin{infobox}[STREAM Algorithm Overview]
\begin{enumerate}
    \item \textbf{Local clustering}: Δημιουργία micro-clusters με BIRCH
    \item \textbf{Weighted k-means}: Clustering των micro-cluster centroids
    \item \textbf{Incremental}: Νέα σημεία ενημερώνουν τα CF χωρίς re-scan
\end{enumerate}
Πλεονεκτήματα: $O(1)$ memory per cluster, single-pass, incremental updates
\end{infobox}

% -----------------------------------------------------------------------------
% Exercise 6.14
% -----------------------------------------------------------------------------
\section{Άσκηση 6.14: Lossy Counting Algorithm}

\begin{formulabox}[Εκφώνηση]
Ο αλγόριθμος Lossy Counting βρίσκει frequent items σε data stream με
error bound $\varepsilon$.

Δίνεται: $\varepsilon = 0.1$, άρα bucket width $w = \lceil 1/\varepsilon \rceil = 10$.

Stream: A, B, A, C, A, B, D, A, C, A, B, A, C, D, A, B, C, A, D, B

\begin{enumerate}
    \item Εκτελέστε τον αλγόριθμο για τα πρώτα 20 στοιχεία
    \item Ποια items έχουν frequency $\geq 0.2$ (support threshold $s = 0.2$)?
\end{enumerate}
\end{formulabox}

\subsection*{Λύση}

\textbf{Lossy Counting Algorithm:}

Δομή δεδομένων: $(item, f, \Delta)$ όπου:
\begin{itemize}
    \item $f$ = estimated frequency
    \item $\Delta$ = maximum possible error (bucket number when item was inserted)
\end{itemize}

\textbf{Κανόνες:}
\begin{itemize}
    \item Νέο item: εισαγωγή με $(item, 1, b_{current} - 1)$
    \item Υπάρχον item: $f \leftarrow f + 1$
    \item Στο τέλος κάθε bucket: διαγραφή αν $f + \Delta \leq b_{current}$
\end{itemize}

\textbf{Εκτέλεση:}

\textbf{Bucket 1} (items 1-10): A, B, A, C, A, B, D, A, C, A

\begin{center}
\begin{tabular}{c|ccc}
\toprule
Item & f & $\Delta$ & $f + \Delta$ \\
\midrule
A & 5 & 0 & 5 \\
B & 2 & 0 & 2 \\
C & 2 & 0 & 2 \\
D & 1 & 0 & 1 \\
\bottomrule
\end{tabular}
\end{center}

End of bucket 1 ($b=1$): Διαγραφή αν $f + \Delta \leq 1$

D: $1 + 0 = 1 \leq 1$ $\to$ \textbf{Διαγράφεται}

\textbf{Bucket 2} (items 11-20): B, A, C, D, A, B, C, A, D, B

\begin{center}
\begin{tabular}{c|ccc}
\toprule
Item & f & $\Delta$ & $f + \Delta$ \\
\midrule
A & 5+3=8 & 0 & 8 \\
B & 2+3=5 & 0 & 5 \\
C & 2+2=4 & 0 & 4 \\
D & 2 & 1 & 3 \\
\bottomrule
\end{tabular}
\end{center}

End of bucket 2 ($b=2$): Διαγραφή αν $f + \Delta \leq 2$

Κανένα item δεν διαγράφεται.

\textbf{(2) Frequent Items με $s = 0.2$:}

Threshold: $sN - \varepsilon N = 0.2 \times 20 - 0.1 \times 20 = 4 - 2 = 2$

Items με $f \geq (s - \varepsilon)N = 2$:

\begin{center}
\begin{tabular}{cccc}
\toprule
Item & $f$ & True freq & $f \geq 2$? \\
\midrule
A & 8 & 8 (40\%) & \checkmark \\
B & 5 & 5 (25\%) & \checkmark \\
C & 4 & 4 (20\%) & \checkmark \\
D & 2 & 3 (15\%) & \checkmark \\
\bottomrule
\end{tabular}
\end{center}

\textbf{Απάντηση:} Items A, B, C έχουν true frequency $\geq 20\%$.

Το D έχει 15\% αλλά λόγω error εμφανίζεται ως candidate (no false negatives, possible false positives).

\begin{warningbox}[Lossy Counting Guarantees]
\begin{itemize}
    \item \textbf{No false negatives}: Όλα τα items με $f \geq sN$ θα βρεθούν
    \item \textbf{Bounded false positives}: Output items έχουν $f \geq (s-\varepsilon)N$
    \item \textbf{Memory}: $O(1/\varepsilon \cdot \log(\varepsilon N))$ entries
\end{itemize}
\end{warningbox}
